\caption{\textbf{Brain-wide stimulus-locked latency tiling, validated on held-out trials; the sparse ``sequence'' cells are mostly unreliable.}
\textbf{a} Feature vectors of 54,569 neurons (odd-numbered trials), sorted by a Rastermap fit to these same trials (Methods parameters with \texttt{grid\_upsample=10}). Display (identical in \textbf{a}, \textbf{b}): per neuron, scaled between median and 99th percentile, values below the 80th percentile white, gamma 0.5. Brackets: dense latency tiling (rows 11,800--16,300) and sparse cells (rows 29,600--35,600; the former ``sequence cells'').
\textbf{b} Even-numbered (held-out) trials in the order of \textbf{a}. Dense tiling and the stimulus-, movement- and feedback-locked blocks replicate; the sharp diagonals of the sparse cells in \textbf{a}, mostly in rarely sampled conditions ($\mathrm{change_b,s}$: $\sim$12 trials per half), do not.
\textbf{c, d, g--j} Neurons are selected on the two concordant stimulus-aligned trial types (odd vs even trials) and evaluated on the four other stimulus-aligned types (even trials), so selection and evaluation share no trials.
\textbf{c} Distribution of each neuron's odd/even reliability (Pearson $r$ between its odd- and even-trial responses; dotted lines, medians): sparse cells are the least reliable (median $r=0.14$; dense tiling 0.81; other 0.23).
\textbf{d} Peak-latency replication on held-out types (Spearman $\rho$) per rate octile: at matched rate, sparse cells replicate less than other cells (the cluster collects the least reliable cells of each rate band).
\textbf{e, f} Zooms of \textbf{b} (even trials, Rastermap order fitted on odd trials, same display): the 6,000 neurons of the sparse cells (\textbf{e}, rows 29,600--35,600) and the 4,500 of dense tiling (\textbf{f}, rows 11,800--16,300), stimulus-aligned trial types. The latency tiling fitted on odd trials persists on even trials for dense tiling, but is not visible for the sparse cells, whose small reproducible part is only revealed by selecting its reliable cells (\textbf{g--i}).
\textbf{g} Peak latency on held-out vs selection types for reliable cells ($r\geq0.3$; sparse $n=1{,}644$, $\rho=0.47$; dense $n=3{,}689$, $\rho=0.81$).
\textbf{h} Correlation of response patterns between fast and slow reaction-time trials (even trials, split at each session's median), aligned to the stimulus (filled) or to movement onset (hatched): both groups are stimulus-locked.
\textbf{i} Latency distributions of the reliable cells overlap.
\textbf{j} Fraction of reliable cells ($r\geq0.3$): sparse cells 32\%, dense tiling 87\%, all other cells 43\%.
\textbf{k, l} Dense tiling cells with added noise matched, neuron by neuron, to the sparse cells' reliability distribution (noise scaled by $1/\sqrt{\text{trials}}$ per trial type), sorted by their own Rastermap fit on odd trials; \textbf{k} odd, \textbf{l} even trials.
\textbf{m} Noise alone turns dense tiling into sparse-cell statistics: median reliability, fraction of reliable cells, cross-validated time--time correlation at lag 0 (odd vs even bins) and peak-latency replication ($\rho$).
Rates in feature-vector units; latencies within the 0--150\,ms stimulus window.}
\label{fig:latency_tiling}
